\documentclass[10pt]{article}

% graphicx package, useful for including eps, pdf, and raster graphics
\usepackage{graphicx}
\DeclareGraphicsExtensions{.png,.png,.jpg,.pdf}

% basic packages
\usepackage{fancyhdr}
\usepackage{color}
\usepackage{parskip}
\usepackage{float}
\usepackage{microtype}
\usepackage{url}
\usepackage{hyperref}
\usepackage{booktabs}
\usepackage{makecell}
\usepackage{multirow}
\usepackage{pbox}
\usepackage{enumitem}
\usepackage{xeCJK}

% fonts
\usepackage{fontspec}
\setmainfont{DroidSans}[
  Extension   = .ttf,
  UprightFont = *,
  BoldFont    = *-Bold ]
\setmonofont{DroidSansMono}[
  Extension   = .ttf,
  UprightFont = * ]
\setCJKmainfont{NotoSansCJK}[
  Extension   = .ttc,
  UprightFont = *-Regular,
  BoldFont    = *-Bold,
  Language    = Chinese Simplified ]

\hypersetup{colorlinks,urlcolor=blue}
\urlstyle{same}

% page layout
\usepackage[top=1in, bottom=1in, left=0.75in, right=0.75in, includefoot, heightrounded]{geometry}

% Useful for seeing the layout as you tweak it.
%\usepackage{showframe}

\setlength{\headheight}{1.1in}
\pagestyle{fancy}
\fancyhf{}

\lhead{\includegraphics[height=1in,clip]{logo}}
\rhead{\pbox[b][1in][c]{\textwidth}{
    \small
    \textbf{检测执行方：}\\
    Northwest Genomics Center\\
    University of Washington School of Medicine\\
    Genome Sciences Box 355065\\
    3720 15th Ave NE\\
    Seattle, WA 98195-5065}}

\newcommand{\PageLine}{\rule{\textwidth}{0.25mm}}

\begin{document}

\begin{description}[font=\normalfont,align=left,labelwidth=6em]
\item [\textbf{\VAR{pat_name|e}}]
\item [\textbf{出生日期：}] \VAR{birth_date|e}
\item [\textbf{标本：}] 鼻腔中部，自我采集
\item [\textbf{采集条形码：}] \VAR{qrcode|e}
\item [\textbf{采集日期：}] \VAR{collect_ts|e}
\item [\textbf{报告日期：}]
  %- if status_code == "never-tested"
  没有进行检测
  %- else
  \VAR{result_ts|e}
  %- endif
\end{description}

\PageLine

\begin{center}
\Large
\textbf{研究结果 – 鼓励您与您的医疗保健提供者共享此结果}
\end{center}

\bigskip

\large \underline{\textbf{关于本检测}}

\begin{itemize}
\item

  COVID-19 是一种名为 SARS-CoV-2 的\textbf{病毒}引发的\textbf{疾病}。
  本研究检测在从您的鼻子里采集的样本中寻找这种病毒，以确定您是否患有 COVID-19。
  这种检测的技术术语叫``逆转录聚合酶链反应''，简称 ``RT-PCR''。

\end{itemize}

\bigskip

\large \underline{\textbf{重要提示}}

查看您的研究检测结果前，您应该查看以下内容：

\begin{itemize}
\item

  您收集的\textbf{样本类型}和\textbf{实验室检测}是\textbf{研究}的一部分内容，
  本报告介绍了您的\textbf{研究检测结果。}

\item

  您自己（或为他人）采集样本并贴上标签后，通过快递将其送回实验室。
  它不是由护士或医生等有执照的医疗保健专业人员采集、贴标签或运送的。这可能会影响您的检测结果，
  医疗保健提供者可能希望再次对您进行检测。

\item

  参与本\textbf{研究}并不能代替常规医疗护理。如果您感觉很不舒服或担心您的检测结果，
  您应该联系医生或医疗保健提供者。

\end{itemize}

\bigskip
\PageLine

\large \underline{\textbf{您的 SARS-CoV-2 (COVID-19) 研究检测结果}}

%- if status_code == "negative"
未检出（阴性）\\
%- elif status_code == "positive"
检出（阳性）\\
%- elif status_code == "inconclusive"
不确定\\
%- elif status_code == "never-tested"
无法进行检测\\
%- else
\VAR{0/0} %# poor man's raise
%- endif

\PageLine
\bigskip

\large \underline{\textbf{了解您的结果}}

此表给出了您的检测可能显示的所有结果，并描述了每个结果的含义。
您可以打印或通过电子邮件将此\textbf{研究检测结果}与您的医生或医疗保健提供者共享。

不管您的检测结果如何，如果您感觉不舒服，您应该待在家里，避免和其他人在一起。
\textbf{如果您一直有症状、您病得很重或担心您自己的疾病，或对您的研究检测结果有疑问，
您应该联系医生或医疗保健提供者。}

\begin{itemize}


\item

  \textbf{If your test says ``Not Detected (Negative)''}, it means our lab did not find the
  virus that causes COVID-19 in the sample that you sent to us. No test is
  perfect, and there is a small chance that you could have the virus even though
  your test is negative.

\item

  \textbf{If your test says ``Detected (Positive)''}, it means our
  lab found the virus that causes COVID-19 in the sample that you sent us. You
  should stay home and avoid being around other people, especially if you are
  feeling sick.

  Your result will be shared with the Washington State and local health
  departments. You do \textbf{not} need to report these results yourself. The
  health department may call you to ask questions about your illness and people
  you may have been around while you were sick.

\item

  \textbf{If your test says ``Inconclusive''}, it means our lab is not certain
  of your test result. No test is perfect, and an inconclusive result happens
  when we find a small amount of SARS-CoV-2 in your sample, but not enough for
  us to say that you have COVID-19.

\item

  \textbf{If your test says ``Unable to be tested''}, it means there was either
  not enough material on your sample for us to test, or there was an error when
  we ran the test. Since we are not able to test your sample, we could not find
  whether or not you have the virus that causes COVID-19.

\end{itemize}

\bigskip

\large \underline{\textbf{资源}}

如果您想了解更多关于 COVID-19 的信息，请查看公共卫生合作伙伴提供的这些资源:

\begin{itemize}
\item
  Public Health \textemdash Seattle \& King County（西雅图和金县公共卫生）\\
  \url{https://www.kingcounty.gov/depts/health/covid-19/languages/chinese.aspx}

\item

  Washington State Department of Health（华盛顿州卫生部）\\
  \url{https://www.doh.wa.gov/Emergencies/NovelCoronavirusOutbreak2020COVID19/ChineseSimplified}

\item
  美国 Centers for Disease Control and Prevention（CDC，疾病控制和预防中心）\\
  \url{https://chinese.cdc.gov/coronavirus/2019-ncov/index.html?CDC\_AA\_refVal=https%3A%2F%2Fwww.cdc.gov%2Fcoronavirus%2F2019-ncov%2Findex-Chinese.html}

\end{itemize}

如果您想通过电话咨询，请拨打以下其中一个电话：

\begin{itemize}

\item

  \textbf{与 COVID-19 有关的医疗问题：} 早上 8 点至晚上 7 点，可拨打
  \textbf{206-477-3977} 联系 King County Call Center（金县呼叫中心）

\item

  \textbf{关于 COVID-19 的非医学问题，包括合规性和企业相关问题:} 周一至周五早上 8:30
  至下午 4:30，可拨打 \textbf{206-296-1608} 联系 King County COVID-19 Business and
  Community Information Line（金县 COVID-19 企业和社区信息热线）

\item

  \textbf{关于华盛顿州 COVID-19 的一般问题:} 早上 6 点至晚上 10 点，可拨打
  \textbf{1-800-525-0127} 联系 Washington State Novel Coronavirus Call Center
  （华盛顿州新型冠状病毒呼叫中心）

\end{itemize}

\bigskip

\large \underline{\textbf{方法和局限性}}

\textbf{方法}\\
本实验使用逆转录 PCR 检测 SARs-CoV-2。该实验使用病毒片段 ORf1B 和 S 的扩增，还检测人类
RnaseP 扩增对照。该检测由 University of Washington 的 Northwest Genomics Center
开发并确定其性能特征。该实验室获得 WA Department of Health医学检测站的许可，
具备进行高复杂性临床实验室检测的资格。

\textbf{局限性}

\begin{itemize}

\item

  这些结果只是研究的一部分内容。

\item

  您自己（或为他人）采集样本并贴上标签后，通过快递将其送回实验室。
  它不是由护士或医生等医疗保健专业人员采集、贴标签或运送的。这可能会影响您的检测结果，
  医疗保健提供者可能希望再次对您进行检测。

\item

  由于此报告基于一项研究数据，因此您的结果可能无法诊断为由 SARS-CoV-2 感染引起的 COVID-19
  疾病。

\end{itemize}

\end{document}
